\documentclass{beamer}
\usepackage[utf8]{inputenc}
\usepackage{listings}
\usepackage{xcolor}

\lstset{
    language=SQL,
    basicstyle=\ttfamily\small,
    keywordstyle=\color{blue},
    commentstyle=\color{gray},
    breaklines=true
}

\title{Modern SQL}
\subtitle{15-445/645 FALL 2024}
\author{Prof. Andy Pavlo}
\institute{Carnegie Mellon University}
\date{}

\begin{document}

\frame{\titlepage}

\begin{frame}{Today's Agenda}
    \begin{itemize}
        \item Database Systems Background [cite: 301]
        \item Relational Model [cite: 301]
        \item Relational Algebra [cite: 301]
        \item Alternative Data Models [cite: 301]
        \item Q\&A Session [cite: 301]
    \end{itemize}
    \vfill
    \hfill 2
\end{frame}

\begin{frame}{Last Class}
    \begin{itemize}
        \item We introduced the Relational Model as the superior data model for databases. [cite: 301]
        \item We then showed how Relational Algebra is the building blocks that will allow us to query and modify a relational database. [cite: 302]
    \end{itemize}
    \vfill
    \hfill 3
\end{frame}

\begin{frame}{SQL History}
    \begin{itemize}
        \item In 1971, IBM created its first relational query language called SQUARE. [cite: 303]
        \item IBM then created "SEQUEL" in 1972 for IBM System R prototype DBMS (Structured English Query Language). [cite: 304, 305]
        \item IBM releases commercial SQL-based DBMSs:
        \begin{itemize}
            \item System/38 (1979) [cite: 305]
            \item SQL/DS (1981) [cite: 305]
            \item DB2 (1983) [cite: 305]
        \end{itemize}
    \end{itemize}
    \vfill
    \hfill 4
\end{frame}

\begin{frame}{SQL Standards Evolution}
    \begin{itemize}
        \item ANSI Standard in 1986. ISO in 1987 (Structured Query Language). [cite: 307]
        \item Current standard is \textbf{SQL:2023}. [cite: 307]
        \begin{itemize}
            \item SQL:2023 $\rightarrow$ Property Graph Queries, Multi-Dim. Arrays [cite: 307]
            \item SQL:2016 $\rightarrow$ JSON, Polymorphic tables [cite: 307]
            \item SQL:2011 $\rightarrow$ Temporal DBs, Pipelined DML [cite: 307]
            \item SQL:2008 $\rightarrow$ Truncation, Fancy Sorting [cite: 307]
            \item SQL:2003 $\rightarrow$ XML, Windows, Sequences, Auto-Gen IDs [cite: 307]
            \item SQL:1999 $\rightarrow$ Regex, Triggers, OO [cite: 308]
        \end{itemize}
        \item The minimum language syntax for SQL support is \textbf{SQL-92}. [cite: 308]
    \end{itemize}
    \vfill
    \hfill 6
\end{frame}

\begin{frame}{Relational Languages}
    \begin{itemize}
        \item Data Manipulation Language (DML) [cite: 327]
        \item Data Definition Language (DDL) [cite: 327]
        \item Data Control Language (DCL) [cite: 327]
        \item Also includes: View definition, Integrity \& Referential Constraints, Transactions. [cite: 327]
        \item \textbf{Important:} SQL is based on bags (duplicates), not sets (no duplicates). [cite: 327]
    \end{itemize}
    \vfill
    \hfill 10
\end{frame}

\begin{frame}{Aggregates}
    Functions that return a single value from a bag of tuples: [cite: 329]
    \begin{itemize}
        \item \texttt{AVG(col)}: Return average value [cite: 329]
        \item \texttt{MIN(col)}: Return minimum value [cite: 330]
        \item \texttt{MAX(col)}: Return maximum value [cite: 330]
        \item \texttt{SUM(col)}: Return sum of values [cite: 330]
        \item \texttt{COUNT(col)}: Return number of values [cite: 331]
    \end{itemize}
    \vfill
    \hfill 13
\end{frame}

\begin{frame}[fragile]{Aggregate Example}
    Get \# of students with a "@cs" login: [cite: 332]
    \begin{lstlisting}
SELECT COUNT(login) AS cnt
FROM student 
WHERE login LIKE '%@cs'
    \end{lstlisting}
    Also equivalent: \texttt{COUNT(*)}, \texttt{COUNT(1)}, \texttt{COUNT(1+1+1)}. [cite: 333, 334, 335]
    \vfill
    \hfill 17
\end{frame}

\begin{frame}[fragile]{Group By}
    Project tuples into subsets and calculate aggregates against each subset. [cite: 342]
    \begin{lstlisting}
SELECT AVG(s.gpa), e.cid
FROM enrolled AS e JOIN student AS s
  ON e.sid = s.sid
GROUP BY e.cid
    \end{lstlisting}
    \textbf{Note:} Non-aggregated values in SELECT must appear in GROUP BY. [cite: 343]
    \vfill
    \hfill 22
\end{frame}

\begin{frame}[fragile]{Having}
    Filters results based on aggregation computation (like a WHERE clause for GROUP BY). [cite: 346, 347]
    \begin{lstlisting}
SELECT AVG(s.gpa) AS avg_gpa, e.cid
FROM enrolled AS e, student AS s
WHERE e.sid = s.sid
GROUP BY e.cid
HAVING avg_gpa > 3.9;
    \end{lstlisting}
    \vfill
    \hfill 28
\end{frame}

\begin{frame}{Window Functions}
    Performs calculation across tuples related to the current tuple without collapsing them. Supports running totals, ranks, and moving averages. [cite: 374]
    \begin{itemize}
        \item \texttt{SELECT FUNC-NAME(...) OVER (...) FROM tableName} [cite: 376]
        \item Special functions: \texttt{ROW\_NUMBER()}, \texttt{RANK()}. [cite: 378]
        \item \texttt{OVER} specifies grouping (\texttt{PARTITION BY}) and sorting (\texttt{ORDER BY}). [cite: 380, 381, 383]
    \end{itemize}
    \vfill
    \hfill 43
\end{frame}

\begin{frame}[fragile]{Common Table Expressions (CTE)}
    Specify a temporary result set referenced by another part of the query (like a temp table for one query). [cite: 413, 414]
    \begin{lstlisting}
WITH cteSource (maxId) AS (
  SELECT MAX(sid) FROM enrolled
)
SELECT name FROM student, cteSource
WHERE student.sid = cteSource.maxId
    \end{lstlisting}
    \vfill
    \hfill 75
\end{frame}

\begin{frame}{Conclusion}
    \begin{itemize}
        \item SQL is a hot language. [cite: 423]
        \item Lots of NL2SQL tools, but writing SQL is not going away. [cite: 424]
        \item You should (almost) always strive to compute your answer as a single SQL statement. [cite: 425]
    \end{itemize}
    \vfill
    \hfill 77
\end{frame}

\end{document}